% Source: https://artofproblemsolving.com/community/c3h3180131_circles_math_problem
% Stable identity: aops-topic-3180131-q01
Let the unknown centre be $(u,v)$ and radius $R>0$. Passing through $(0,5)$ gives $u^2+(v-5)^2=R^2$. For each given circle, write the squared centre distance as $(R+s)^2$, with $s=1$ for external tangency and $s=-1$ for internal tangency. The unknown circle cannot lie inside either unit circle because $(0,5)$ lies outside both, so an internally tangent solution must have $R>1$. For the two signs $s,t$, subtracting the point equation from the two tangency equations gives
\[-8u+10v-10=2sR,\qquad 8u+10v-10=2tR.\]
If $s=t=1$, then $u=0$, $v=1+\dfrac{R}{5}$, and the positive-radius solution is $R=\dfrac{10}{3},v=\dfrac{5}{3}$. If $s=t=-1$, the solution is $R=5,v=0$. If $s=-t$, then $v=1$ and $u=\pm\dfrac{R}{4}$, so $R^2=\dfrac{R^2}{16}+16$, giving $R=\dfrac{16}{\sqrt{15}}$ and $u=\pm\dfrac{4}{\sqrt{15}}$. All four satisfy the required tangencies and pass through the point, so the answer is $\boxed4$.
\begin{asy}
usepackage("amsmath");
size(7cm);
for(int k=0;k<4;++k){pair off=(14*k,0),O;real r;if(k==0){O=off+(0,5/3);r=10/3;}else if(k==1){O=off;r=5;}else{O=off+((k==2?1:-1)*4/sqrt(15),1);r=16/sqrt(15);}draw(circle(off+(-4,0),1));draw(circle(off+(4,0),1));draw(circle(O,r),blue);dot(off+(0,5));label("$P$",off+(0,5),N);dot(O);label("$O$",O,S);label("Case "+string(k+1),off+(0,-6));}
\end{asy}
